The rapid evolution of Visual Language Models (\acsp{VLM}) has progressively expanded the ability of neural systems to jointly process perceptual and linguistic information. This line of research builds on the hypothesis that linguistic meaning cannot be constructed exclusively from distributional regularities learned from textual corpora, but must also be grounded in perceptual experience \cite{baroni2015grounding}. Within this context, Visual Question Answering (\acs{VQA}) emerged as one of the first systematic paradigms for assessing whether a model can answer questions about the content of an image \cite{antol2015vqa}. Subsequent benchmarks, such as GQA, extended this setting by introducing compositional questions involving entities, attributes, and visual relations \cite{hudson2019gqa}.
However, high task accuracy alone does not provide sufficient evidence of multimodal comprehension. Several studies have shown that models can exploit statistical regularities in answers, questions, or dataset distributions, producing correct predictions without necessarily \textit{understanding} the associated visual content. Evaluation has therefore progressively moved beyond task performance towards the assessment of grounding, understood as the ability to associate a prediction with the relevant perceptual evidence.
One widely adopted strategy for evaluating grounding relies on minimally contrastive pairs, in which a correct instance is paired with a foil obtained by modifying a semantically relevant element. \textit{FOIL-it!}, for example, evaluates whether a model can detect localized inconsistencies between an image and its linguistic description \cite{shekhar2017foilit}. \textit{Winoground}, instead, presents pairs of images and descriptions containing the same lexical elements but arranged according to different compositional relations, showing that recognizing individual entities does not necessarily entail understanding the relational structure connecting them \cite{thrush2022winoground}.
The transition from images to videos introduces additional challenges. Video understanding requires integrating information distributed over time, recognizing actions and state changes, ordering events, and establishing temporal and causal relations. \textit{Action Genome Question Answering (\acs{AGQA})} addresses these requirements through compositional questions concerning objects, actions, and spatio-temporal relations \cite{grundemclaughlin2021agqa}. The \textit{Multi-modal Video Understanding Benchmark (\acs{MVBench})} similarly evaluates a broader range of competencies, including action perception, movement, temporal ordering, and dynamic reasoning \cite{li2024mvbench}.
Nevertheless, video benchmarks remain vulnerable to unimodal shortcuts and to strategies that rely predominantly on the most informative frame. \textit{Video Language Model Assessment (\acs{VILMA})} addresses this limitation as a zero-shot benchmark based on video--text contrastive pairs \cite{kesen2024vilma}. For each complex capability, \acs{VILMA} introduces a proficiency test designed to verify the preliminary perceptual ability required to solve the corresponding task. The results show that, once these prerequisite conditions are considered, part of the apparently correct performance can be attributed to accidental or spurious strategies. In particular, the examined \acsp{VLM} do not display a systematic advantage over statistical image-based models on tasks requiring strict temporal grounding.
Within this line of research, \acs{MAIA} introduced a natively Italian benchmark designed to provide a fine-grained evaluation of video-linguistic reasoning \cite{testa2025maia}. The resource covers twelve reasoning categories related to linguistic phenomena, including causal, counterfactual, implicit, spatial, temporal, sentimental, planning, uncertainty, and out-of-scope reasoning. Its main contribution lies in aligning two tasks built on the same videos and questions: \textit{Visual Statement Verification (\acs{VSV})}, where the model selects the correct statement from a minimally contrastive pair, and \textit{Open-Ended Visual Question Answering (\acs{OEVQA})}, where the model freely generates the answer.
This setting makes it possible to assess not only accuracy but also coherence between discriminative comprehension and linguistic generation. A model may correctly answer a multiple-choice question while failing to reproduce the same information in an open-ended setting, or may behave inconsistently across equivalent linguistic formulations. To capture this phenomenon, \acs{MAIA} introduces the \textit{Aggregate Accuracy} metric, which rewards only those instances in which the model correctly answers all true--false pairs associated with the same question and simultaneously generates the correct response in the corresponding open-ended task. The reduction in Aggregate Accuracy relative to conventional single-instance accuracy highlights the fragility of apparently strong performance and the need to jointly evaluate robustness, consistency, and grounding.
Although \acsp{VLM} have traditionally focused on text--vision interactions, \underline{real videos are intrinsically audiovisual}. Their interpretation may depend on music, prosody, background noise, and other acoustic signals that cannot be reconstructed from visual frames alone. An increasing body of research has accordingly extended multimodal language model architectures to support joint processing of text, images, video, and audio.
\textit{ImageBind}, for instance, introduced a shared representation space spanning images, text, audio, depth, thermal information, and inertial movement \cite{girdhar2023imagebind}. Other multimodal language models (\acsp{MLM}) rely on modality-specific encoders combined with projection or fusion modules that align heterogeneous representations with the linguistic space. \textit{Video-SALMONN: Speech-Enhanced Audio-Visual Large Language Models} introduces a multi-resolution Q-former\footnote{A multi-resolution Q-former, such as a Multi-Resolution Causal Q-Former, processes continuous audio, video, or image information at different temporal or spatial scales through distinct query banks, balancing fine-grained signals with broader contextual information.} to align audio signals with synchronized video \cite{sun2024videosalmonn}, while \textit{VideoLLaMA2} combines spatio-temporal modeling and acoustic processing within a generative multimodal architecture \cite{cheng2024videollama2}.
More recent omnimodal architectures process multiple modalities within a unified framework. \textit{Qwen2.5-Omni} adopts a thinker--talker architecture that accepts text, image, video, and audio inputs and produces textual or audio responses \cite{xu2025qwen25omni}. The subsequent \textit{Qwen3.5-Omni} architecture integrates visual and audio encoders into a common backbone and introduces explicit temporal references to improve audio--video alignment \cite{qwen2026qwen35omni}.
These developments demonstrate that heterogeneous signals can be processed within a single architecture. The availability of a modality, however, does not imply that it is used functionally or complementarily. The distinction between \textbf{multimodal availability} and \textbf{multimodal integration} is essential here: the former indicates that a system can accept multiple input sources, the latter requires the model to select, coordinate, and combine relevant information across modalities. A system may thus be multimodal from an architectural perspective while relying predominantly on unimodal information during prediction.
This issue is evident in \textit{Audio-Visual Question Answering} benchmarks such as MUSIC-\acs{AVQA}, which evaluates spatio-temporal relations in musical performance videos and extends the analysis to objects, entities, and everyday activities \cite{li2022musicavqa}. Subsequent studies reveal a strong visual bias, showing that many questions can be answered using the video stream alone, without exploiting the audio channel. Under these conditions, similar audio-only and video-only performance does not demonstrate genuine integration, but rather suggests a form of unimodal collapse in which one modality dominates the decision process.
To address this problem, \textit{\acs{DAVE}} was introduced as a benchmark in which each question jointly requires visual and acoustic information, making either modality insufficient in isolation \cite{radevski2026dave}. The benchmark distinguishes among action recognition, sound classification, temporal ordering, and multimodal synchronization, a decomposition that helps determine whether an error arises from the failure to perceive an event or from the inability to establish temporal relations between sounds and visual actions. The results show that models may achieve strong performance on individual components while encountering greater difficulty in audio--visual synchronization, indicating that cross-modal integration remains a central limitation.
The increasing attention to audio has also led to benchmarks specifically designed to assess acoustic intelligence. \textit{\acs{MMAUPro}} evaluates 49 competencies distributed across speech, environmental sounds, music, and their combinations \cite{kumar2026mmaupro}. These tasks include long-audio comprehension, multi-track reasoning, spatial localization, prosody interpretation, and the integration of dialogue, music, and background noise. The persistently limited performance of current models indicates that acoustic processing cannot be reduced to dialogue recognition alone, since audio also conveys spatial, temporal, emotional, and contextual information that must be evaluated independently.
A complementary perspective is provided by \textit{\acs{SONIC01}}, a benchmark for omnimodal models based on audio--visual interactions drawn from real conversational and everyday contexts \cite{radwan2026sonic}. \acs{SONIC01} combines open-ended summarization, multiple-choice questions, and temporal event localization, while also showing that replacing raw audio with transcription removes paralinguistic information related to tone and communicative intention. Its results identify temporal localization as one of the most challenging tasks, confirming that the integration of verbal content, acoustic signals, and visual sequences remains an open problem.
The comparison between native audio and transcription is also central to my recent work, \textit{Lost in Translation: Investigating the Role of Audio Information for Video-based Visual Statement Verification} \cite{deodati2026lost}. The study extends the \acs{MAIA} framework by introducing audio into the \acs{VSV} task and comparing two representations of the acoustic modality: raw audio directly processed by \textit{Qwen2.5-Omni-7B} and automatic transcriptions supplied as additional textual input to \textit{Qwen2.5-VL-7B}.
To characterize the audio content, the videos are divided into four classes:
\begin{description}
\item[Music:] Instances characterized by repetitive vocalizations, onomatopoeic patterns, strongly repetitive tokens, or explicit lexical cues associated with non-speech audio.
\item[Dialogue:] Videos whose transcriptions display sufficiently coherent linguistic content, acceptable log-probability values, low degeneracy, and limited evidence of no-speech segments.
\item[Silence/Noise:] Videos with empty or highly unreliable transcriptions, particularly when associated with low log-probability values and relatively high no-speech probabilities.
\item[Unknown:] Residual cases that appear speech-related but do not satisfy the dialogue criteria with sufficient confidence, thereby preserving potentially informative content while explicitly marking its uncertainty.
\end{description}
Class labels and speech transcriptions were obtained by combining the outputs of \textit{Whisper-1} and \textit{\acs{GPT}-4o-transcribe}, together with an assessment of output plausibility. This classification distinguishes results according to the informativeness of the available audio content.
The results reveal a marked dominance of the visual channel. Audio provides only a limited benefit, mainly when no visual information is available. Once rich visual input is provided, audio's contribution becomes negligible and may even introduce noise or errors, with limited robustness observed when a single black frame is added instead. Under visually rich conditions, the difference between with-audio and without-audio settings is minimal, suggesting that audio functions more as a surrogate channel than as information fully integrated with vision.
Audio transcriptions exhibit a different behavior. When added to already rich visual input, they reduce accuracy relative to the visual-only condition. This shows that transcription is not a neutral transformation: once projected into the same representational space as the questions and answer options, it can compete with them, increasing the influence of textual regularities and introducing redundant or irrelevant information.
\textit{Lost in Translation} thus shows that different representations of the same acoustic content can affect the model's decision process differently. Raw audio can be partially ignored when it is not useful, whereas transcription interacts directly with the linguistic component of the model. Introducing an additional modality does not therefore guarantee an improvement, and may instead compromise predictions that were already correct on the basis of reliable visual evidence.
These findings must be interpreted in light of the original \acs{MAIA} design. The benchmark was created to investigate visual content, and annotators were explicitly instructed not to consider audio information, so its questions and answers were not designed to exploit acoustic evidence. The limited contribution of audio may therefore reflect both the models' modality-integration mechanisms and the limited relevance of audio within the original dataset.
Recent literature further emphasizes the importance of identifying the distinct sources of multimodal errors. \textit{\acs{MIRAGE}} disentangles hallucinations caused by perceptual failures from those arising from incorrect reasoning over correctly represented information \cite{dong2025mirage}. Although primarily focused on visual reasoning, its methodological principle extends naturally to audio--visual settings, where an incorrect response may derive from failed sound recognition, inaccurate temporal synchronization, selection of the wrong modality, or erroneous inference over otherwise correctly represented evidence.
Complementary attribution approaches seek to identify which parts of the input contribute to the generation of an answer. \textit{OmniTrace} proposes an attribution framework for omnimodal autoregressive generation that links generated tokens to textual segments, visual regions, and audio/video intervals, aggregating attention signals into cross-modal explanations associated with linguistically interpretable units. This type of analysis is particularly relevant for determining whether a model actually exploits audio information or whether its prediction is driven primarily by text and vision.
The current state of the art thus reveals a clear tension between the rapid development of omnimodal architectures and their effective ability to integrate heterogeneous modalities. Contemporary models can receive multiple input types within the same context, yet diagnostic benchmarks provide no guarantee that these sources are used complementarily. Performance may instead depend on a dominant modality, linguistic shortcuts, dataset correlations, or additional signals that introduce noise.
A central unresolved limitation is the construction of benchmarks in which audio and visual information are explicitly supervised. Assessing genuine audio--visual comprehension requires items in which audio contributes relevant information and the correct answer cannot be determined from a single modality alone. Such a design would distinguish complementarity, dominance, substitution, and interference, while evaluating whether the model selects the appropriate evidence and combines it coherently across modalities.
The present study follows this direction. Building on the competence-oriented \acs{MAIA} framework and the recent findings of \textit{Lost in Translation}, it investigates model behavior on a dataset specifically designed to require a genuine contribution from audio in order to answer the questions correctly. The central hypothesis is that, when audio and video provide partial and complementary evidence, multimodal comprehension depends on the model's ability to reconstruct the missing information through cross-modal integration, avoiding both unimodal collapse and linguistic shortcuts based on plausibility.

